% TOP_PROBLEM: {"id": 1101307, "title": "Questions in Geometric Group Theory — Q 13.7", "category_id": 17, "category": {"id": 17, "name": "group_theory", "display_name": "Group Theory", "description": "Problems about groups, group actions, representations, and related algebraic structures.", "slug": "group-theory", "order_index": 17, "created_at": "2026-07-31T15:26:25.671Z"}, "status": "open", "problem_number": "AMR-010-1307", "set_id": 13}
% FIRST_POSTED: 2026-10-01
% ENTRY_KIND: statement_only
% UnsolvedMath ID 1101307; source catalogue code AMR-010-1307
% !TeX program = lualatex
% Definition and sources only; this file is not an LLM solution attempt.
\documentclass[11pt,a4paper]{article}
\usepackage[margin=25mm]{geometry}
\usepackage{fontspec}
\setmainfont{lmroman10-regular.otf}[BoldFont=lmroman10-bold.otf,
 ItalicFont=lmroman10-italic.otf,BoldItalicFont=lmroman10-bolditalic.otf]
\usepackage{amsmath,amssymb,mathtools}
\usepackage{unicode-math}
\setmathfont{latinmodern-math.otf}
\AtBeginDocument{\let\setminus\smallsetminus}
\usepackage{xurl,hyperref}
\hypersetup{colorlinks=true,urlcolor=blue}
\setcounter{secnumdepth}{0}
\setlength{\parindent}{0pt}
\setlength{\parskip}{5pt}
\setlength{\emergencystretch}{3em}
\begin{document}
\section*{Questions in Geometric Group Theory — Q 13.7}\label{1101307}
{\small UnsolvedMath ID 1101307 · AMR-010-1307 · Group Theory · AMR Open Problem Lists}
\subsection{Definitions and mathematical statement}
Let $k\geq1$ and let $L$ be a finite simplicial complex whose underlying topological space is homeomorphic to the sphere $S^{2k-1}$. The complex is flag if every finite set of vertices which are pairwise joined by edges spans a simplex. This condition rules out a collection of mutually adjacent vertices whose expected simplex is missing.

For $0\leq i\leq2k-1$, let $f_i(L)$ be the number of $i$-dimensional simplices, counted without orientations. Define the rational number
\[
 C(L)=1-\sum_{i=0}^{2k-1}(-1)^i\frac{f_i(L)}{2^{i+1}}.
\]
The Charney--Davis question asks whether every such flag triangulation satisfies
\[
 (-1)^k C(L)\geq0.
\]
The initial constant $1$ is part of this expression; $C(L)$ is not simply the ordinary Euler characteristic of the sphere. The exponent $k$ is half of one more than the sphere dimension, exactly as displayed.

Both the integer $k$ and the finite triangulation vary universally. No assumption that the sphere is the boundary of a convex polytope is made, and no particular flag subdivision is prescribed. A proof for polytopal spheres or low dimensions would therefore cover only a restricted family. The assertion is non-strict, so equality is permitted. All the quantities depend on the simplicial face numbers, and the problem asks for a numerical inequality valid for every permissible combinatorial realization of the sphere, rather than invariance of this number under arbitrary subdivisions.
\subsection{Short English statement}
Does every flag triangulation of an odd-dimensional sphere satisfy the signed Charney–Davis face-number inequality?
\subsection{Sources}
\begin{itemize}
\item[S1] ULAM.AI and the UnsolvedMath Contributors, supplied problems.json, record 1101307 (AMR-010-1307), AMR Open Problem Lists. Catalogue identity and source transcription; CC BY 4.0.
\url{https://huggingface.co/datasets/ulamai/UnsolvedMath}.
\item[S2] Bestvina - Questions in Geometric Group Theory (pdf) (2004), Question 13.7 (PDF page 25). Original source indexed by AMR-010-1307.
\url{https://www.math.utah.edu/~bestvina/eprints/questions-updated.pdf}.
\end{itemize}
\end{document}
